\begin{textsample}{NEG Dim 5 – vem\_ed\_29 – Score 1.00 – vem\_ed\_29/t151.txt}  \label{ex:f5_neg_126}
Leitores Aos Osvaldo Succi Jr.
Coordenador PCIs Os Projetos Colaborativos Internacionais (PCIs) são parte da área técnica de Apoio à Internacionalização do Ensino Superior, na Divisão de Extensão e Pesquisa do Ensino Superior (Depes) da Unidade do Ensino Superior de Graduação (Cesu) do Centro Paula Souza.
Em 8 de maio de 2025, a equipe organizou a quarta edição do International Virtual Exchange Symposium (IVES Cesu 2025), evento online dedicado a debater a situação dos Intercâmbios Virtuais a partir de diversas perspectivas internacionais.
Neste ano, contamos com as palestras de João Monteiro, do ISP Gaya (Portugal), Stephanie Doscher, da University of Minnesota (EUA), Julia Bacca e Diana Garzon, da Uniminuto (Colômbia) e Adam Freed, da University of Michigan (EUA).
Uma síntese do evento encontra-se nesta edição, que também relata minha apresentação na Conferência Internacional Faubai 2025, realizada presencialmente em Brasília, de 26 a 30 de \textbf{abril}.
Na seção “Artigo de Opinião”, temos o relato de Neusa Haruka Sezaki Gritti, responsável pelo apoio administrativo e pedagógico à equipe de Apoio à Internacionalização do Ensino Superior e professora de Inglês nas Fatecs Mogi das Cruzes e Sebrae, sobre a importância da troca cultural na realização do PCI.
Quer contribuir você também?
Acesse https://publicacoescesu.cps.sp.gov.br/VEm/about/submissions Dúvidas?
Escreva para cesu.pci@cps.sp.gov.br.
Boa leitura!

% matched lemmas: abril
\end{textsample}
